%!TeX root=MemoriaTFG.tex

\chapter{Estado del arte}
\label{cap:sota}

\section{Aprendizaje profundo aplicado a dermatología}
\label{sec:dl-dermatologia}

El trabajo de Esteva et~al.~\cite{esteva2017} establece la primera
referencia consolidada de clasificación dermatológica con redes
convolucionales, con rendimiento equivalente al de un panel de
dermatólogos certificados sobre clasificación binaria
maligno/benigno. Las extensiones posteriores amplían el espectro
diagnóstico y la diversidad de datos~\cite{tschandl2018,ham10000} y
formulan tareas adicionales como segmentación binaria de
lesión~\cite{isic2018} y clasificación binaria
melanoma/no-melanoma con anotaciones estructuradas
multi-criterio~\cite{kawahara2019}. El paradigma común permanece
estable: entrenamiento desde pesos de ImageNet, una arquitectura por
tarea y dependencia de anotación experta. Tres limitaciones lo
acotan: el coste de la anotación clínica, la fragmentación
taxonómica entre conjuntos públicos y la pérdida de rendimiento
fuera del dominio de entrenamiento, particularmente acusada en
fototipos cutáneos elevados~\cite{fitzpatrick17k,ddi}.

\section{Modelos fundacionales: arquitectura y objetivos}
\label{sec:fundacionales-base}

Los modelos fundacionales sustituyen el entrenamiento específico por
encoders preentrenados de forma auto-supervisada o multimodal sobre
datasets visuales masivos. La arquitectura predominante es el Vision
Transformer~\cite{vit2021}, que procesa la imagen como una secuencia
de parches y aprende dependencias globales mediante mecanismos de
atención. Los objetivos de preentrenamiento más extendidos son tres:
el \emph{masked image modeling}, en el que la red reconstruye
parches ocultos a partir del contexto visible; la destilación
auto-supervisada entre vistas; y la alineación contrastiva
imagen-texto al estilo CLIP~\cite{clip2021}. La adaptación posterior
a tareas específicas opera por sondeo lineal sobre representaciones
congeladas, ajuste fino supervisado, o transferencia directa
zero-shot guiada por \emph{prompts} en los modelos con rama
lingüística.

\section{Encoders visuales auto-supervisados}
\label{sec:encoders-visuales}

DINOv2~\cite{dinov2} es la referencia abierta de encoder visual
auto-supervisado de propósito general, entrenado sobre
$\sim 142$ millones de imágenes naturales sin etiquetas mediante
destilación entre vistas. Sus representaciones transfieren a tareas
de clasificación, recuperación y segmentación sobre dominios
naturales y biomédicos con ajuste mínimo, y constituyen una base
habitual para adaptación a dominios especializados.

En el dominio dermatológico, PanDerm~\cite{panderm2025} es un
encoder visual ViT (Base $\sim 86$\,M y Large $\sim 307$\,M
parámetros) entrenado sobre $\sim 2{,}1$ millones de imágenes
dermatológicas en cuatro modalidades (fotografía corporal total,
dermatoscopia, fotografía clínica e histopatología). El
preentrenamiento utiliza \emph{masked latent modeling}: la red
reconstruye representaciones latentes de regiones ocultas en lugar
de píxeles, con un modelo CLIP-Large congelado como
\emph{teacher}. Sus autores reportan una mejora media del
$+10{,}2\%$ sobre el estado del arte previo en $28$ tareas
dermatológicas. La disponibilidad simultánea de pesos, código y
benchmarks lo convierte en el primer modelo fundacional dermatológico
verificable de forma independiente.

\section{Modelos vision-lenguaje contrastivos}
\label{sec:vision-lenguaje}

CLIP~\cite{clip2021} establece el preentrenamiento contrastivo
imagen-texto sobre $400$ millones de pares de la web. El objetivo
InfoNCE aproxima la imagen a su descripción dentro del lote y la
aleja del resto. El espacio compartido habilita dos modos
operativos: clasificación zero-shot guiada por \emph{prompts}
textuales (la clase ganadora es la de mayor similitud coseno entre
el vector visual y los vectores textuales de las clases candidatas)
y recuperación bidireccional imagen-texto sobre datasets
precomputados.

SigLIP~\cite{siglip2023} sustituye la pérdida softmax de CLIP por
una pérdida sigmoide que reduce la dependencia del tamaño de batch y
escala con mejor estabilidad. Su variante Large
(SO400M) dispone de embedding de $1\,152$ dimensiones y
$\sim 878$\,M parámetros. BiomedCLIP~\cite{biomedclip} extiende
el paradigma al dominio biomédico sobre PMC-15M
($\sim 15$ millones de pares figura-leyenda de literatura
científica); su cobertura dermatológica específica es limitada por
el sesgo de selección del dataset.

En dermatología, Derm1M y la familia DermLIP~\cite{derm1m} añaden la
rama lingüística contrastiva al estado del arte de la disciplina.
Derm1M reúne $1\,029\,761$ pares imagen-texto sobre
$403\,563$ imágenes únicas, organizados según una ontología
jerárquica de cuatro niveles construida por dermatólogos. La familia
DermLIP combina varios encoders visuales (PanDerm Base,
ConvNeXt-Large) con varios encoders textuales (PubMedBERT extendido
a $256$ tokens, GPT-2 dermatológico con tokenizador de CLIP a
$77$ tokens). La configuración con mejor rendimiento reportada
por los autores combina PanDerm Base como encoder visual y
PubMedBERT como encoder textual, superando a BiomedCLIP en los
benchmarks zero-shot evaluados.

La implementación operativa de la recuperación visual densa sobre
datasets a gran escala requiere índices vectoriales optimizados para
búsqueda aproximada del vecino más próximo. FAISS~\cite{faiss2019}
y estructuras basadas en grafos del tipo Hierarchical Navigable
Small Worlds permiten consultas en el orden de $10^2$
milisegundos sobre datasets de centenares de miles a millones de
\emph{embeddings} precomputados.

\section{Modelos integrados con zero-shot nativo}
\label{sec:zero-shot-nativo}

DermFM-Zero~\cite{dermfmzero2026} integra el preentrenamiento visual
auto-supervisado y la alineación contrastiva imagen-texto en un
único modelo con capacidad zero-shot nativa. El entrenamiento se
estructura en dos fases: \emph{masked latent modeling} sobre tres
millones de imágenes en la primera, y \emph{bootstrapped contrastive
learning} sobre un millón de pares texto-imagen en la segunda. El
encoder textual es PubMedBERT. El modelo incorpora interpretabilidad
emergente mediante Sparse Autoencoders~\cite{templeton2024}, redes
auxiliares que descomponen la representación interna en un
diccionario disperso de direcciones asociables a conceptos
clínicamente identificables. Sus autores reportan validación en
tres \emph{reader studies} con más de mil cien clínicos. Sus pesos
no son públicos en la fecha de cierre del presente trabajo, lo que
restringe su tratamiento a la discusión conceptual.

BLIP-2~\cite{blip2_2023} introduce un patrón arquitectónico
distinto: conecta un encoder visual congelado con un modelo de
lenguaje preentrenado mediante un módulo intermedio ligero
(\emph{Q-Former} o \emph{Querying Transformer}) que transforma las
representaciones visuales en una secuencia reducida de tokens
interpretables por el modelo de lenguaje. Sólo el Q-Former se
entrena. No existe, en la fecha de cierre, una adaptación pública
específica de BLIP-2 al dominio dermatológico.

\section{Segmentación promptable}
\label{sec:segmentacion-sota}

SAM~\cite{sam2023} establece el paradigma de segmentación
promptable: produce máscaras a partir de \emph{prompts} visuales
(puntos, cajas, máscaras parciales) sobre imagen natural. Su
preentrenamiento utiliza $\sim 11$ millones de imágenes con
más de mil millones de máscaras generadas mediante un proceso
semi-automatizado. SAM2~\cite{sam2} amplía el enfoque al dominio
del vídeo y mejora el rendimiento sobre imágenes estáticas.

La transferencia a dominios médicos requiere adaptación supervisada
sobre conjuntos con máscaras binarias. La estrategia habitual
combina un modelo SAM/SAM2 preentrenado con \emph{fine-tuning} sobre
ISIC2018~\cite{isic2018} y adaptadores LoRA~\cite{lora}, que
incorporan un número reducido de parámetros entrenables manteniendo
congelado el peso original del encoder visual. Esta estrategia
reduce los requisitos de memoria y cómputo y permite la adaptación
sobre hardware modesto.

\section{Modelos de lenguaje multimodales generalistas}
\label{sec:llm-multimodal}

Los modelos de lenguaje multimodales preentrenados sobre datasets web
masivos constituyen una referencia de comparación inevitable: GPT-4o
(OpenAI), Gemini~2.5~Pro (Google DeepMind) y la familia MedGemma
(versiones 4B Vision y 27B Text, Google). Su número de parámetros
supera al de los encoders fundacionales específicos en uno o dos
órdenes de magnitud. Su modo de uso característico es la inferencia
guiada por \emph{prompts} sin \emph{fine-tuning} específico por
tarea. Trabajos recientes documentan su uso conversacional para
soporte al razonamiento clínico~\cite{amie2024,deepmind_aicoclinician}.
Los modelos de acceso restringido (GPT-4o, Gemini) sólo se evalúan
mediante API; MedGemma libera pesos y se evalúa en local.

\section{Problemas metodológicos y guías de reporte}
\label{sec:problemas-metodologicos}

La aplicación de modelos fundacionales a imagen médica plantea tres
problemas metodológicos directamente relevantes para el presente
trabajo. El primero es la fuga de información entre conjuntos de
evaluación y dataset de preentrenamiento. La construcción de Derm1M
sobre fuentes web abiertas implica que datasets como Dermnet, HIBA
o MSKCC pueden solaparse con su dataset de entrenamiento, lo que
sesga al alza las métricas zero-shot reportadas. El segundo es la
heterogeneidad taxonómica entre datasets, que impide comparaciones
directas sin un paso explícito de armonización. El tercero es la
disparidad de rendimiento por fototipo cutáneo, documentada en
Fitzpatrick17k~\cite{fitzpatrick17k} y atendida específicamente por
SkinCon~\cite{skincon} mediante anotación basada en $48$
conceptos clínicos visuales.

Las guías de reporte recientes para inteligencia artificial en
imagen médica ---CLAIM~\cite{claim2024} y la actualización
TRIPOD+AI--- establecen la transparencia metodológica como
condición de defensibilidad: cohorte, particiones, métricas,
incertidumbre y limitaciones deben constar explícitamente en el
informe completo. El reporte de pruebas estadísticas (intervalos de
confianza por \emph{bootstrap}, comparaciones múltiples,
contrastes pareados de AUROC mediante el test de
DeLong~\cite{delong1988}, y tamaños muestrales por subgrupo)
constituye un requisito creciente en publicaciones de la disciplina.

\section{Brechas operativas y posicionamiento}
\label{sec:brecha}

Cinco brechas operativas delimitan el espacio en el que el presente
trabajo aporta evidencia empírica. \textbf{(i)~Disponibilidad real
de los modelos.} PanDerm y DermLIP liberan pesos y código;
DermFM-Zero no lo hace. BiomedCLIP, DINOv2, SigLIP, SAM, SAM2 y
BLIP-2 son abiertos. Los modelos generalistas ofrecen acceso parcial
por API. Cualquier evaluación empírica externa queda condicionada
por esta disponibilidad. \textbf{(ii)~Benchmarks no homogéneos entre
familias.} Los benchmarks reportados por los autores se construyen
sobre subconjuntos derivados del propio preentrenamiento y aplican
protocolos no equivalentes entre familias. La comparación bajo un
mismo protocolo y un mismo conjunto de datasets externos no está
documentada. \textbf{(iii)~Heterogeneidad taxonómica.} La
comparación inter-dataset exige un esquema de armonización que la
literatura no proporciona de forma consolidada.
\textbf{(iv)~Concentración geográfica y lingüística.} Los datasets
proceden mayoritariamente de instituciones de Australia, Estados
Unidos y Europa central, y los encoders textuales se preentrenan
sobre datasets en inglés. La disponibilidad de material clínico en
español con texto narrativo asociado es marginal.
\textbf{(v)~Distancia entre evaluación experimental e integración
hospitalaria.} La evaluación sobre \emph{benchmarks} públicos no
cubre las condiciones operativas asociadas al despliegue en un
entorno sanitario real (hardware no estándar, restricciones de
privacidad, trazabilidad, interacción con sistemas hospitalarios).

La contribución empírica del trabajo, cuyo diseño se describe en el
capítulo~\ref{cap:metodos} y cuyos resultados se presentan
cuantitativamente en el capítulo~\ref{cap:resultados}, se sitúa
sobre estas cinco brechas.
